\subsection{Pembahasan Costmap}

Setiap lapisan \emph{costmap} memiliki parameter yang dapat diatur secara independen sesuai kebutuhan tiap \emph{region}. \emph{Global costmap} dikonfigurasi melalui nilai \emph{Cost Scaling Factor} (CSF) yang memengaruhi gradien biaya inflasi di sekitar \emph{obstacle}. Semakin tinggi CSF, \emph{path} yang dihasilkan semakin pendek karena robot diperbolehkan melintas lebih dekat ke \emph{obstacle}. Namun, nilai CSF yang terlalu tinggi dapat menyebabkan \emph{path} tidak aman jika \emph{clearance} terhadap \emph{obstacle} terlalu kecil. Untuk mengakomodasi perbedaan karakteristik medan, CSF dapat diatur berbeda pada tiap \emph{region}, misalnya CSF tinggi pada area terbuka dan CSF rendah pada lorong sempit.

\emph{Local costmap} dapat dikonfigurasi melalui  (\emph{window size}) dan resolusi grid. Ukuran jendela yang lebih besar memberikan informasi lingkungan yang lebih luas bagi DWA \emph{local planner} dalam memilih \emph{trajectory}, tetapi memerlukan lebih banyak sumber daya komputasi. Sebaliknya, ukuran jendela yang terlalu kecil dapat menyebabkan DWA tidak memiliki cukup informasi untuk menghindari \emph{obstacle} yang berada di luar jangkauan. Parameter ini perlu disesuaikan dengan kecepatan robot dan kerapatan \emph{obstacle} di lingkungan sekitar.

\emph{Transition costmap} memiliki parameter radius ekspansi dan nilai \emph{cost} pada zona transisi. Radius ekspansi menentukan seberapa jauh zona buffer menyebar dari garis transisi, yang memengaruhi seberapa awal robot mendeteksi bahwa ia akan berpindah \emph{region}. Nilai \emph{cost} pada zona transisi ($E = 5$) dan garis transisi ($T = 1$) digunakan oleh \emph{region switcher} untuk menentukan kapan robot harus bersiap melakukan lompatan dan kapan lompatan dieksekusi.

Kemampuan untuk mengonfigurasi setiap lapisan \emph{costmap} secara terpisah memberikan fleksibilitas dalam menyesuaikan perilaku navigasi terhadap kondisi lingkungan yang berbeda. Pendekatan ini sejalan dengan prinsip \emph{layered costmap}, di mana setiap lapisan bertanggung jawab atas aspek lingkungan yang spesifik dan dapat dioptimalkan tanpa memengaruhi lapisan lainnya.

\subsection{Pembahasan Lokalisasi}

Sistem lokalisasi merupakan hasil kerja sama antara ICP dan \emph{region switcher}. ICP bertanggung jawab melacak posisi robot secara
presisi pada tingkat sub-meter, sedangkan \emph{region switcher} melacak lokasi
robot pada tingkat yang lebih besar, yaitu menentukan \emph{region} mana yang
sedang ditempati robot. Pembagian peran ini memungkinkan sistem untuk tetap
mengetahui posisi robot secara global meskipun ICP mengalami penurunan akurasi
di area tertentu, seperti zona transisi antar \emph{region}.

Proses pelompatan \emph{region} menggunakan mekanisme \emph{transition mode}
yang memanfaatkan data \emph{pitch} dari sensor kompas CMPS14 untuk mendeteksi
perubahan kemiringan saat robot memasuki \emph{slope} pada Region B. Ketika
nilai \emph{pitch} melebihi ambang batas ($12^\circ$), sistem langsung memicu
lompatan tanpa harus menunggu \emph{timer} selesai. Mekanisme ini lebih responsif
dibandingkan \emph{timer-based} karena mendeteksi transisi secara langsung
berdasarkan perubahan fisik lingkungan, bukan berdasarkan estimasi waktu.
\emph{Timer-based transition} tetap dipertahankan sebagai \emph{safety fallback}
jika data \emph{pitch} tidak tersedia atau tidak mencukupi ambang batas.

\subsection{Pembahasan Path Planner}

Dalam sistem navigasi multi-\emph{region}, \emph{global planner} dan \emph{path
router} bekerja sama untuk merencanakan perjalanan dari posisi robot menuju
tujuan akhir. Bentuk \emph{path} pada setiap segmen dapat diubah melalui parameter
\emph{planner} itu sendiri, tidak hanya dari \emph{costmap}. Parameter
seperti \emph{cost factor}, \emph{neutral cost}, dan \emph{Cost Scaling
Factor} memengaruhi bagaimana \emph{global planner} menyeimbangkan antara
jarak tempuh dan jarak aman terhadap \emph{obstacle}. Setiap kali robot
melompat ke \emph{node} baru, \emph{global planner} membuat \emph{path}
baru untuk segmen selanjutnya, sehingga \emph{path} yang digunakan selalu
sesuai dengan posisi dan \emph{region} robot saat itu.

Hasil pengujian menunjukkan bahwa tidak ada satu nilai CSF yang optimal
untuk semua \emph{region}. Pada Region C yang memiliki ruang
terbuka, CSF 100 menghasilkan \emph{path} terpendek namun deviasinya
terhadap \emph{ground truth} paling tinggi. Sebaliknya, pada Region B
yang berupa lorong, seluruh nilai CSF menghasilkan \emph{path}
yang hampir identik. Temuan ini menunjukkan perlunya strategi
perencanaan jalur yang bersifat spesifik per \emph{region}. Dengan
demikian, setiap segmen dapat menggunakan konfigurasi \emph{global
planner} yang berbeda.

\subsection{Pembahasan DWA Local Planner }

Pergerakan robot di lingkungan simulasi menghadapi beberapa lorong sempit,
terutama di Region C. Robot hexapod yang digunakan cukup kuat untuk tidak
mudah tersangkut, namun radius \emph{footprint} tetap perlu dikonfigurasi
dengan tepat agar dapat menghindari tabrakan dengan dinding. Parameter DWA
memungkinkan robot bergerak secara \emph{omnidirectional}, yang
membantu dalam melewati lorong-lorong sempit karena robot dapat bergerak
menyamping tanpa harus berputar terlebih dahulu.

Konfigurasi DWA yang berbeda pada tiap \emph{region} memungkinkan
\emph{trajectory} yang dihasilkan tetap optimal sesuai karakteristik medan.
Pada Region A yang luas, kecepatan tinggi ($0,08$ m/s) digunakan untuk meminimalkan
waktu tempuh. Pada Region B yang berupa lorong sempit, kecepatan diturunkan
($0,04$ m/s) untuk menjaga stabilitas navigasi pada ruang terbatas. Pada Region C
dengan \emph{obstacle} yang kompleks, jumlah sampel kecepatan ditingkatkan
(18 sampel) dan \emph{sim\_time} diperpanjang ($10$ detik) agar DWA memiliki
lebih banyak kandidat lintasan untuk menemukan \emph{trajectory} yang aman.
Selain itu, \emph{goal bias} pada Region C diturunkan menjadi $24$ (dari $32$)
agar DWA lebih fleksibel dalam memilih lintasan yang menghindari \emph{obstacle}
tanpa terlalu terikat pada arah tujuan. Penyesuaian ini menunjukkan bahwa
konfigurasi DWA spesifik per \emph{region} berkontribusi langsung terhadap
keberhasilan navigasi.

\subsection{Pembahasan Navigasi Penuh dari Region A ke C}

\begin{table}[H]
\centering
\footnotesize
\begin{tabular}{c|c|c|c|c|c|c|c}
\toprule
\thead{Percobaan} & \thead{Waktu A\\(s)} & \thead{Waktu T0\\(s)} & \thead{Waktu B\\(s)} & \thead{Waktu T1\\(s)} & \thead{Waktu C\\(s)} & \thead{Total\\(s)} & \thead{Sukses} \\
\midrule
1 & 13.56 & 12.52 & 5.46 & 7.73 & 214.54 & 233.56 & Ya \\
2 & 14.50 & 14.33 & 6.41 & 7.27 & 165.45 & 186.36 & Ya \\
3 & 14.50 & 16.96 & 7.73 & 6.98 & 151.12 & 173.35 & Ya \\
4 & 15.59 & 11.75 & 7.33 & 7.86 & 115.32 & 138.24 & Tidak\\
5 & 15.23 & 13.60 & 7.10 & 6.50 & 129.48 & 151.81 & Ya \\
6 & 19.00 & 16.76 & 7.85 & 7.97 & 200.13 & 226.97 & Ya \\
7 & 19.00 & 16.00 & 8.28 & 8.24 & 194.00 & 221.28 & Ya \\
8 & 21.39 & 18.74 & 15.67 & 17.69 & 103.95 & 141.01 & Tidak \\
9 & 18.43 & 16.55 & 8.77 & 8.04 & 166.53 & 193.74 & Ya \\
10 & 17.48 & 14.97 & 7.83 & 8.48 & 182.42 & 207.73 & Ya \\
Mean & 16.46 & 15.21 & 7.43 & 7.65 & 175.46 & 199.35 & 80\%\\
\bottomrule
\end{tabular}
\caption{Rangkuman waktu tempuh }
\label{tab:dwa_rangkuman_waktu}
\end{table}

Pengujian navigasi penuh dari region A ke region C dilakukan untuk mengevaluasi
hasil keseluruhan sistem navigasi. Data waktu tempuh dibagi menjadi waktu di
tiap region (A, B, dan C) serta waktu robot di zona transisi T0 (antara A dan
B) dan T1 (antara B dan C). Dari 10 kali percobaan, diperoleh 8 keberhasilan
dan 2 kegagalan. Waktu tempuh terbaik adalah 141.01 detik dan waktu terlama
233.56 detik, yang menunjukkan variasi yang cukup besar dalam performa
navigasi antar percobaan.

Variasi waktu tempuh terbesar terjadi di region C. Hal ini disebabkan oleh
karakteristik \emph{DWA planner} yang harus mensimulasikan lintasan
dengan \emph{unknown obstacle}. Ketika robot mendekati \emph{unknown
obstacle} yang terdeteksi oleh LiDAR saat navigasi berlangsung, DWA
menurunkan kecepatan dan memilih lintasan yang lebih aman, sehingga
memperpanjang waktu tempuh. Sebaliknya, pada region~A dan B dengan
lingkungan yang lebih sederhana dan tanpa \emph{unknown obstacle}, durasi
navigasi relatif konsisten antar percobaan. Hal ini terlihat dari nilai Total
yang bervariasi antara 141.01 hingga 233.56 detik, di mana variasi tersebut
hampir seluruhnya berasal dari fluktuasi waktu di region C (103.95-214.54
detik) akibat perbedaan jumlah dan posisi \emph{unknown obstacle} yang
dihadapi robot pada setiap percobaan.

Dua kegagalan terjadi pada percobaan ke-4 dan ke-8. Pada percobaan ke-4,
robot gagal menghasilkan lintasan setelah keluar dari \emph{transition mode}
di region~C. \emph{DWA planner} tidak menemukan \emph{trajectory} yang valid
dalam \emph{dynamic window} akibat konfigurasi kecepatan dan orientasi robot
yang tidak ideal pasca lompatan, sehingga sistem tidak dapat melanjutkan
navigasi. Pada percobaan ke-8, kegagalan disebabkan oleh hilangnya lokalisasi
ICP setelah melewati \emph{transition mode}. Koréksi pose yang tidak stabil
pasca lompatan menyebabkan \emph{local planner} menerima referensi posisi
yang salah, sehingga navigasi tidak dapat dilanjutkan. Kedua kegagalan ini
menunjukkan bahwa \emph{transition mode} merupakan titik kritis dalam sistem
navigasi multi-\emph{region}.

Secara keseluruhan, seluruh komponen navigasi (lokalisasi ICP, \emph{global
planner} A*, \emph{path router}, \emph{region switcher}, dan \emph{DWA local
planner}) terintegrasi dengan baik dan mampu melakukan navigasi dari region~A
ke region~C secara otonom. Keberhasilan 8 dari 10 percobaan menunjukkan
tingkat keandalan yang memadai, dengan dua kegagalan yang disebabkan oleh
kondisi tepi pada \emph{transition mode} yang dapat diminimalkan melalui
optimalisasi parameter \emph{transition costmap} dan penanganan \emph{edge
case} pada mekanisme lompatan.